\chapter{能量法}

\intro{宁鸣而生，不默而死。}

\section{应变能的计算}

\subsection{外力功与应变能}
弹性体在缓慢加载过程中，外力所作的功 \(W\) 全部转化为弹性应变能 \(U\)，即
\[
U = W.
\]
对于线性弹性体，外力与位移成线性关系。例如轴向拉压杆，力 \(F\) 与伸长量 \(\Delta l\) 满足 \(F = k \Delta l\)，此时外力功为
\[
W = \frac{1}{2} F \Delta l.
\]
因此应变能
\[
U = \frac{1}{2} F \Delta l.
\]

\begin{myTheorem}{基本变形下的应变能公式}
    \begin{itemize}
        \item \textbf{轴向拉压}（等直杆，轴力 \(F_N\) 为常数）：
        \[
        U = \frac{F_N^2 l}{2EA} = \frac{EA}{2l} (\Delta l)^2.
        \]
        \item \textbf{圆轴扭转}（扭矩 \(T\) 为常数）：
        \[
        U = \frac{T^2 l}{2GI_p}.
        \]
        \item \textbf{平面弯曲}（纯弯，弯矩 \(M\) 为常数）：
        \[
        U = \frac{M^2 l}{2EI}.
        \]
    \end{itemize}
\end{myTheorem}

若杆件受组合变形，且各内力分量仅为单个变量的函数，则总应变能等于各内力分量引起应变能的代数和。更一般地，对于长度为 \(l\) 的杆，应变能可写为积分形式：
\begin{myformula}
    U = \int_l \frac{F_N^2}{2EA}\,\mathrm{d}x + \int_l \frac{T^2}{2GI_p}\,\mathrm{d}x + \int_l \frac{M^2}{2EI}\,\mathrm{d}x.
\end{myformula}
这里略去了剪力引起的应变能（通常较小）。

\subsection{克拉比隆定理}
对于线性弹性结构，若有 \(n\) 个广义力 \(F_1,F_2,\dots,F_n\) 作用，相应的广义位移为 \(\delta_1,\delta_2,\dots,\delta_n\)，且载荷按比例加载，则总应变能为
\[
U = \frac{1}{2} \sum_{i=1}^n F_i \delta_i.
\]

\section{卡氏定理}

卡氏定理（Castigliano's Theorem）是能量法中求位移和力的有力工具。

\begin{myTheorem}{卡氏第一定理}
    弹性体的应变能 \(U\) 表示成广义位移 \(\delta_1,\delta_2,\dots,\delta_n\) 的函数时，有
    \[
    F_i = \frac{\partial U}{\partial \delta_i},
    \]
    即广义力等于应变能对相应广义位移的偏导数。
\end{myTheorem}
此定理可用来由位移求力。

\begin{myTheorem}{卡氏第二定理（卡氏位移定理）}
    对于线性弹性结构，若将应变能 \(U\) 表示为载荷的函数，则应变能对任一载荷 \(F_i\) 的偏导数等于该载荷作用点沿载荷方向的位移 \(\delta_i\)：
    \[
    \delta_i = \frac{\partial U}{\partial F_i}.
    \]
    更一般地，若需求与载荷对应的广义位移，同样适用。
\end{myTheorem}

应用卡氏第二定理时，若所求位置无载荷，可附加一虚设载荷 \(F_0\)，计算 \(\partial U/\partial F_0\) 后再令 \(F_0=0\) 即得真实位移。

\begin{example}
    简支梁 \(AB\)，跨长 \(l\)，受集中力 \(F\) 作用于跨度中点。求中点挠度 \(w_C\)。
\end{example}
\begin{solution}
    取坐标系，梁的弯矩方程为分段函数。对于 \(0 \le x \le l/2\)：
    \[
    M(x) = \frac{F}{2}x, \quad \frac{\partial M}{\partial F} = \frac{x}{2}.
    \]
    由对称性，全梁应变能 \(U = 2\int_0^{l/2} \frac{M^2}{2EI}\mathrm{d}x\)。
    根据卡氏第二定理，中点挠度
    \[
    w_C = \frac{\partial U}{\partial F} = 2\int_0^{l/2} \frac{M}{EI} \frac{\partial M}{\partial F}\mathrm{d}x
    = 2\int_0^{l/2} \frac{1}{EI}\left(\frac{F}{2}x\right)\left(\frac{x}{2}\right)\mathrm{d}x
    = \frac{F}{2EI} \int_0^{l/2} x^2 \mathrm{d}x = \frac{Fl^3}{48EI}.
    \]
\end{solution}

\section{莫尔定理（单位载荷法）}

莫尔定理是求结构任意点位移的通用方法，适用于线性与非线性弹性结构。

\begin{myTheorem}{莫尔定理}
    结构在真实载荷作用下，任一点沿某一方向的位移 \(\Delta\) 可由下式计算：
    \[
    \Delta = \int_l \frac{F_N \overline{F}_N}{EA}\,\mathrm{d}x + \int_l \frac{T \overline{T}}{GI_p}\,\mathrm{d}x + \int_l \frac{M \overline{M}}{EI}\,\mathrm{d}x.
    \]
    其中 \(F_N, T, M\) 为真实载荷引起的内力；\(\overline{F}_N, \overline{T}, \overline{M}\) 为在所求位移点沿位移方向作用单位广义力（力或力偶）引起的内力。
\end{myTheorem}

莫尔积分实际上就是卡氏第二定理的推广，特别适用于只需计算某一点位移的情形。

\begin{example}
    悬臂梁 \(AB\)，长 \(l\)，自由端受集中力 \(F\)，求自由端的挠度和转角。
\end{example}
\begin{solution}
    \begin{enumerate}
        \item \textbf{求挠度}：在自由端沿竖直方向加单位力 \(1\)。\\
        采用坐标系：原点在固定端，\(x\) 轴指向自由端。则弯矩
        \[
        M(x) = -F(l-x), \quad \overline{M}(x) = -(l-x).
        \]
        代入莫尔积分（忽略剪力和轴力）：
        \[
        w_B = \int_0^l \frac{M \overline{M}}{EI}\mathrm{d}x = \int_0^l \frac{[-F(l-x)][-(l-x)]}{EI}\mathrm{d}x
        = \frac{F}{EI} \int_0^l (l-x)^2 \mathrm{d}x = \frac{Fl^3}{3EI}.
        \]
        \item \textbf{求转角}：在自由端加单位力偶 \(1\)。\\
        实际弯矩同上 \(M(x) = -F(l-x)\)；单位力偶引起弯矩 \(\overline{M}(x) = -1\)。
        \[
        \theta_B = \int_0^l \frac{M \overline{M}}{EI}\mathrm{d}x = \int_0^l \frac{[-F(l-x)](-1)}{EI}\mathrm{d}x
        = \frac{F}{EI} \int_0^l (l-x)\mathrm{d}x = \frac{Fl^2}{2EI}.
        \]
    \end{enumerate}
\end{solution}

\section{图乘法}

对于等截面直杆，当莫尔积分中的 \(M\) 图与 \(\overline{M}\) 图至少有一个为线性变化时，积分可简化为弯矩图面积与形心对应另一弯矩图竖标的乘积，即图乘法。

\begin{myformula}
    \int_l \frac{M \overline{M}}{EI}\,\mathrm{d}x = \frac{1}{EI} \sum \omega \cdot \overline{M}_C,
\end{myformula}
其中 \(\omega\) 为 \(M\) 图的面积，\(\overline{M}_C\) 为与 \(M\) 图形心对应的 \(\overline{M}\) 图的竖标。

使用规则：
\begin{itemize}
    \item 将 \(M\) 图分段，使每段内 \(\overline{M}\) 图呈线性；
    \item 若 \(\overline{M}\) 图发生转折，须在转折处分段；
    \item 面积 \(\omega\) 与竖标 \(\overline{M}_C\) 同侧时乘积为正，异侧为负。
\end{itemize}

常见图形的面积和形心位置（抛物线、三角形等）应熟记。
\subsection*{图乘法常用图形面积与形心位置}
\begin{table}[htbp]
    \centering
    \caption{常见弯矩图形的面积 \(\omega\) 与形心位置 \(\bar{x}\)（距图形左端）}
    \label{tab:figure-integral}
    \begin{tabular}{@{} l c c c @{}}
        \toprule
        \textbf{图形形状} & \textbf{面积 \(\omega\)} & \textbf{形心位置 \(\bar{x}\)} & \textbf{说明} \\
        \midrule
        矩形 & \(b h\) & \(\dfrac{b}{2}\) & —— \\
        \addlinespace
        三角形（顶点在左端） & \(\dfrac{1}{2} b h\) & \(\dfrac{2}{3} b\) & 直角在左端 \\
        \addlinespace
        三角形（顶点在右端） & \(\dfrac{1}{2} b h\) & \(\dfrac{1}{3} b\) & 直角在右端 \\
        \addlinespace
        二次抛物线（顶点在左端） & \(\dfrac{1}{3} b h\) & \(\dfrac{3}{4} b\) & \(y = kx^2\) 型 \\
        \addlinespace
        二次抛物线（顶点在右端） & \(\dfrac{1}{3} b h\) & \(\dfrac{1}{4} b\) & 倒置抛物线 \\
        \addlinespace
        二次抛物线（对称顶点） & \(\dfrac{2}{3} b h\) & \(\dfrac{b}{2}\) & 如均布荷载弯矩图 \\
        \addlinespace
        三次抛物线（顶点在左端） & \(\dfrac{1}{4} b h\) & \(\dfrac{4}{5} b\) & 较少使用 \\
        \addlinespace
        梯形 & \(\dfrac{b}{2}(h_1+h_2)\) & \(\dfrac{b}{3}\cdot\dfrac{h_1+2h_2}{h_1+h_2}\) & 可拆为矩形+三角形 \\
        \bottomrule
    \end{tabular}
\end{table}





\begin{myhypothesis}{图乘法使用规则}
    \begin{enumerate}
        \item 仅适用于 \textbf{等截面直杆}（\(EI\) 为常数）。
        \item 两个弯矩图中至少有一个是 \textbf{线性变化}（直线或折线）。
        \item 若 \(\overline{M}\) 图有转折，必须在其 \textbf{转折处分段}，分别图乘后求和。
        \item 面积 \(\omega\) 取正值，竖标 \(\overline{M}_C\) 的正负按“\textbf{同侧为正，异侧为负}”确定。
        \item 若 \(M\) 图较复杂，可分解为简单图形（矩形、三角形、抛物线）分别图乘后叠加。
    \end{enumerate}
\end{myhypothesis}

\section{互等定理}

\subsection{功的互等定理（贝蒂-瑞利互等定理）}
\begin{myTheorem}{功的互等定理}
    对于线性弹性结构，第一力系在第二力系引起的位移上所作的功等于第二力系在第一力系引起的位移上所作的功。
    \[
    \sum_{i} F_i^{(1)} \delta_i^{(2)} = \sum_{j} F_j^{(2)} \delta_j^{(1)}.
    \]
\end{myTheorem}

\subsection{位移互等定理（麦克斯韦互等定理）}
\begin{myTheorem}{位移互等定理}
    若在结构上 \(i\) 点作用广义力 \(F_i=1\)，在 \(j\) 点引起的广义位移为 \(\delta_{ji}\)；在 \(j\) 点作用广义力 \(F_j=1\)，在 \(i\) 点引起的广义位移为 \(\delta_{ij}\)，则
    \[
    \delta_{ij} = \delta_{ji}.
    \]
\end{myTheorem}
位移互等定理在结构分析中非常有用，可大大减少计算量。

\begin{example}
    一简支梁，在 \(A\) 端作用力偶 \(M_A\) 时，测得跨度中点 \(C\) 的转角为 \(\theta_C\)。问在 \(C\) 点作用集中力 \(F\) 时，\(A\) 端的转角是多少？
\end{example}
\begin{solution}
    由功的互等定理，力系1：\(M_A\)（力偶）；力系2：\(F\)（集中力）。设 \(\theta_A\) 为力系2引起的 \(A\) 端转角，\(\delta_C\) 为力系1引起的 \(C\) 点挠度。则
    \[
    M_A \cdot \theta_A^{(2)} = F \cdot \delta_C^{(1)}.
    \]
    又已知在 \(M_A\) 作用下 \(C\) 点转角为 \(\theta_C\)，但这里位移不直接对应。若利用位移互等定理，当 \(M_A=1\) 时，\(C\) 点挠度 \(\delta_{CA}\) 与 \(F=1\) 时 \(A\) 端转角 \(\theta_{AC}\) 相等。因此可根据已知条件直接得出 \(\theta_A = \frac{F}{M_A}\delta_C\) 等。
\end{solution}



